% --- 1. INFORMAZIONI GENERALI ---
\section{Informazioni Generali}
\begin{itemize}
	\item \textbf{Azienda Coinvolta:} Zucchetti\textsubscript{G}
	\item \textbf{Mezzo di Comunicazione:} Google Meet
	\item \textbf{Data:} 2026-05-26
	\item \textbf{Orario di inizio:} 12:00
	\item \textbf{Orario di fine:} 12:40
\end{itemize}

Il presente verbale documenta l’incontro, svoltosi tramite Google Meet, 
tra il gruppo Synergo Unit e il referente dell’azienda Zucchetti. 
Durante la riunione sono stati presentati e discussi il PoC(Proof of Concept), i problemi dell'infastruttura, i requisiti non funzionali e la disponibilità dell'azienda Zucchetti durante il periodo estivo.



\vspace{0.5cm}
\textbf{Referenti Azienda:}
\begin{table}[ht]
	\centering
	\renewcommand{\arraystretch}{1.2}
	\begin{tabularx}{\textwidth}{X l}
		\toprule
		\textbf{Nome e Cognome} & \textbf{Ruolo\textsubscript{G}} \\
		\midrule
		Gregorio Piccoli & Referente Principale \\
		\bottomrule
	\end{tabularx}
\end{table}

% --- 2. CHIARIMENTI TECNICI RICHIESTI ---
\section{Chiarimenti Richiesti}
Il gruppo ha richiesto chiarimenti sui seguenti aspetti:

\begin{enumerate}
	\item Proof of Concept
	\item Problemi dell'infastruttura
	\item Requisiti non funzionali
	\item Disponibilità dell'azienda durante il periodo estivo
\end{enumerate}

% --- 3. RESOCONTO ---
\section{Resoconto}

\subsection{Proof of Concept}

Il Proof of Concept è stato visionato e approvato.

\subsection{Problemi dell'infastruttura}

Per quanto riguarda l’infrastruttura attuale, è stato suggerito di utilizzare inizialmente il modello *gemma4:12b*. Eventualmente, potranno essere ampliati i nostri permessi per consentire l’utilizzo del modello *Qwen3.6-35B-A3B*, ritenuto più stabile e con prestazioni superiori. È stato inoltre raccomandato di evitare chiamate a cascata tra modelli, in quanto l’infrastruttura prevede il ricaricamento del modello ad ogni richiesta, con conseguente perdita di efficienza.


\subsection{Requisiti non funzionali}
I requisiti non funzionali dovranno essere ampliati e resi maggiormente misurabili. In particolare:
\begin{itemize}
    \item \textbf{Compatibilità browser:} è necessario definire i browser ufficialmente supportati e le eventuali versioni minime. Browser come Edge, Safari e Chrome condividono la stessa base tecnologica e non dovrebbero presentare criticità rilevanti; al contrario, Firefox utilizza un motore differente e richiede test specifici, così come Safari su macOS.
    \item \textbf{Prestazioni AI:} occorre chiarire quale tempo di attesa sia considerato accettabile per la generazione di una proposta AI. In precedenza era stato indicato un tempo minimo di 2 secondi; si è deciso tuttavia di non fissare un limite massimo, introducendo invece un meccanismo di interruzione manuale che consenta all’utente di interrompere l’elaborazione.
    \item \textbf{Persistenza opzionale:} è stato chiarito che il requisito opzionale relativo al database comprende funzionalità complete di autenticazione e gestione delle note (visualizzazione elenco, apertura, aggiornamento ed eliminazione), e non la sola memorizzazione. L'eliminazione delle note è considerata meno rilevante in quanto rappresenta una funzionalità standard che non introduce elementi di particolare innovazione.
\end{itemize}

\subsection{Disponibilità dell'azienda durante il periodo estivo}

Per quanto riguarda il periodo estivo, considerata la presenza di ferie, il docente ha suggerito di valutare la possibilità di individuare preventivamente degli slot di comunicazione, al fine di evitare ritardi dovuti a difficoltà di coordinamento. È stato inoltre comunicato che l’azienda resterà chiusa nelle due settimane centrali di agosto e che il referente sarà assente per l’intero mese. Si raccomanda pertanto di concentrare le comunicazioni entro la fine di luglio e l’inizio di settembre, mantenendo come eventuali finestre operative la prima e l’ultima settimana di agosto in caso di necessità.
\newpage
% --- 4. DECISIONI FORMALI ---
\section{Decisioni Formali}
Nella seguente tabella sono tracciate le decisioni adottate a seguito
dei chiarimenti ricevuti.

\renewcommand{\arraystretch}{1.3}
\vspace{-0.3cm}
\noindent
\begin{longtable}{c p{12cm}}
	\toprule
	\textbf{Codice} & \textbf{Decisione} \\
	\midrule
	\endfirsthead

	\toprule
	\textbf{Codice} & \textbf{Decisione} \\
	\midrule
	\endhead

	\midrule
	\multicolumn{2}{r}{\textit{Continua nella pagina successiva}} \\
	\midrule
	\endfoot

	\bottomrule
	\endlastfoot

	\textbf{VE-7.1}
		& Proof of Concept approvato \\[4pt]

	\textbf{VE-7.2}
		& Scelte riguardo i problemi dell'infastruttura \\[4pt]

	\textbf{VE-7.3}
		& Scelte riguardo i requisiti funzionali \\[4pt]

	\textbf{VE-7.4}
		& Discussione della disponibilità aziendale nel periodo estivo \\[4pt]
\end{longtable}

\newpage
% --- 5. AZIONI (TODO) ---
\section{Azioni da Svolgere (To-Do)}

Le seguenti attività sono state individuate come conseguenza
delle decisioni adottate durante la riunione.

\renewcommand{\arraystretch}{1.3}
\begin{longtable}{c c p{5cm} p{3cm} l}
	\toprule
	\textbf{Codice} & \textbf{Rif. Decisione} & \textbf{Descrizione Task} & \textbf{Assegnatario} & \textbf{Scadenza} \\
	\midrule
	\endfirsthead

	\toprule
	\textbf{Codice} & \textbf{Rif. Decisione} & \textbf{Descrizione Task} & \textbf{Assegnatario} & \textbf{Scadenza} \\
	\midrule
	\endhead

	\midrule
	\multicolumn{5}{r}{\textit{Continua nella pagina successiva}} \\
	\midrule
	\endfoot

	\bottomrule
	\endlastfoot

\end{longtable}

% --- 6. FIRME ---
\section*{Approvazione e Firme}
Con le seguenti firme, i rappresentanti approvano formalmente il contenuto
del presente verbale e le decisioni adottate.

\vspace{2.5cm}

\noindent
\begin{minipage}[c]{0.45\textwidth}
	\begin{center}
		\rule{6cm}{0.4pt} \\
		\vspace{0.2cm}
		\textbf{Per il Gruppo Synergo Unit} \\
		Armando Moda Scarati \\
		\textit{Responsabile\textsubscript{G} di Progetto}
	\end{center}
\end{minipage}
\hfill
\begin{minipage}[c]{0.45\textwidth}
	\begin{center}
		\rule{6cm}{0.4pt} \\
		\vspace{0.2cm}
		\textbf{Per l'Azienda Zucchetti} \\
		Gregorio Piccoli \\
		\textit{Referente Aziendale}
	\end{center}
\end{minipage}

\vspace{2cm}